% assets/w09-lidar-mount.svg — 선수추진기 WAM-V 에 3D LiDAR 를 어디에 다는가 (실측, 2026-10-27)
%
%   bash _tools/tikz2svg.sh assets/src/w09-lidar-mount.tex assets/w09-lidar-mount.svg
%
% 이 그림이 말하는 것 하나
%   라이다 높이는 계산이 아니라 **실측**으로 정했다. 1.6 m 로 계산해 두었더니
%   들어온 값의 82 % 가 장애물이 아니라 자기 배였다. 윗면이 z=1.765 m 인 전장함
%   (cpu_cases_link) 때문이다. 1.80 부터 0 % 가 되고, 파도에 기우는 것까지 보아 1.90 으로 둔다.
%
% 선체 좌표계 : x 전방, y 좌현, z 상방 (URDF/ROS 관례 — FLU)
% 화면 : x 오른쪽, z 위쪽.  1 m = 15 mm
% 측면도의 7.09 m 는 그림이 너무 길어져 치수선을 끊어 그렸다 (중간의 빗금 두 개)

\documentclass[border=6pt]{standalone}
% capstone-style.tex — 이 볼트의 TikZ 그림이 공유하는 스타일
%
% 쓰는 법:  % capstone-style.tex — 이 볼트의 TikZ 그림이 공유하는 스타일
%
% 쓰는 법:  % capstone-style.tex — 이 볼트의 TikZ 그림이 공유하는 스타일
%
% 쓰는 법:  \input{capstone-style}   (assets/src/ 안에서)
% 빌드:     bash _tools/tikz2svg.sh assets/src/w02-node-topic.tex assets/w02-node-topic.svg
%
% 왜 TikZ 인가
%   손으로 SVG 좌표를 쓰면 화살촉·선 굵기·글자 크기가 그림마다 미묘하게 달라진다.
%   그것이 "자료가 어설퍼 보인다"의 정체다. 스타일을 한 번만 정하고 상대좌표로
%   배치하면 좌표를 고쳐도 그림이 무너지지 않는다.
%   수식은 본문과 같은 TeX 글꼴로 나온다.
%
% 한글
%   ko.TeX(DVI 경로)를 쓴다. XeLaTeX 은 TikZ 그래픽을 PDF special 로 내보내는데
%   dvisvgm 이 그것을 받으려면 Ghostscript < 10.01 이 필요하다. 이 PC 에는
%   10.07 이 깔려 있어 그 경로가 막힌다. latex -> DVI -> dvisvgm 이 유일하게
%   동작하는 조합이다.

\usepackage[hangul]{kotex}
\usepackage{tikz}
\usepackage{amsmath}
\usepackage{xcolor}
\usetikzlibrary{arrows.meta, positioning, calc, fit, backgrounds, decorations.pathmorphing}

% ---- 색 --------------------------------------------------------------------
% 강조는 하나만. 나머지는 검정. 흑백 인쇄에서도 읽혀야 한다.
\definecolor{ink}{HTML}{1A1A1A}
\definecolor{soft}{HTML}{666666}
\definecolor{accent}{HTML}{7C3AED}
\definecolor{warn}{HTML}{B91C1C}

% 계층 색 — Simulink 모델의 블록 색과 같은 규칙을 쓴다
\definecolor{guid}{HTML}{1D4ED8}    % 유도
\definecolor{ctrl}{HTML}{EA580C}    % 제어
\definecolor{thr}{HTML}{D97706}     % 추진기
\definecolor{plant}{HTML}{16A34A}   % 운동모델
\definecolor{nav}{HTML}{7C3AED}     % 항법 · 신호처리
\definecolor{pathc}{HTML}{0D9488}   % 계획 경로 · 항적
\definecolor{errc}{HTML}{C2410C}    % 오차
\definecolor{flow}{HTML}{0369A1}    % 조류 · 바람

% ---- 선과 화살촉 -----------------------------------------------------------
\tikzset{
  sig/.style   = {ink, line width=0.5pt, -{Stealth[length=4pt, width=3pt]}},
  wire/.style  = {ink, line width=0.5pt},
  fb/.style    = {ink, line width=0.5pt, densely dashed,
                  -{Stealth[length=4pt, width=3pt]}},
  asig/.style  = {accent, line width=0.5pt, densely dashed,
                  -{Stealth[length=4pt, width=3pt]}},
  % 블록 — 흰 바탕, 직각, 검은 테두리. 색은 테두리에만 준다
  blk/.style   = {draw=ink, line width=0.55pt, fill=white,
                  minimum width=18mm, minimum height=8.5mm, inner sep=2pt, font=\small},
  gblk/.style  = {blk, draw=guid},
  cblk/.style  = {blk, draw=ctrl},
  tblk/.style  = {blk, draw=thr},
  pblk/.style  = {blk, draw=plant},
  nblk/.style  = {blk, draw=nav},
  % 합산점 · 분기점
  sum/.style   = {draw=ink, line width=0.55pt, fill=white, circle,
                  inner sep=0pt, minimum size=4.6mm, font=\scriptsize},
  dot/.style   = {ink, fill=ink, circle, inner sep=0pt, minimum size=1.5mm},
  % 글자
  sname/.style = {font=\small, inner sep=1.5pt},
  note/.style  = {font=\scriptsize, soft, inner sep=1.5pt},
  wnote/.style = {font=\scriptsize, warn, inner sep=1.5pt},
  anote/.style = {font=\scriptsize, accent, inner sep=1.5pt},
  axis/.style  = {ink, line width=0.5pt},
  tick/.style  = {font=\scriptsize, soft},
  % 묶음 상자 — 서브시스템 표시
  grp/.style   = {draw=soft, line width=0.4pt, densely dotted, rounded corners=1mm,
                  inner sep=3mm},
}
   (assets/src/ 안에서)
% 빌드:     bash _tools/tikz2svg.sh assets/src/w02-node-topic.tex assets/w02-node-topic.svg
%
% 왜 TikZ 인가
%   손으로 SVG 좌표를 쓰면 화살촉·선 굵기·글자 크기가 그림마다 미묘하게 달라진다.
%   그것이 "자료가 어설퍼 보인다"의 정체다. 스타일을 한 번만 정하고 상대좌표로
%   배치하면 좌표를 고쳐도 그림이 무너지지 않는다.
%   수식은 본문과 같은 TeX 글꼴로 나온다.
%
% 한글
%   ko.TeX(DVI 경로)를 쓴다. XeLaTeX 은 TikZ 그래픽을 PDF special 로 내보내는데
%   dvisvgm 이 그것을 받으려면 Ghostscript < 10.01 이 필요하다. 이 PC 에는
%   10.07 이 깔려 있어 그 경로가 막힌다. latex -> DVI -> dvisvgm 이 유일하게
%   동작하는 조합이다.

\usepackage[hangul]{kotex}
\usepackage{tikz}
\usepackage{amsmath}
\usepackage{xcolor}
\usetikzlibrary{arrows.meta, positioning, calc, fit, backgrounds, decorations.pathmorphing}

% ---- 색 --------------------------------------------------------------------
% 강조는 하나만. 나머지는 검정. 흑백 인쇄에서도 읽혀야 한다.
\definecolor{ink}{HTML}{1A1A1A}
\definecolor{soft}{HTML}{666666}
\definecolor{accent}{HTML}{7C3AED}
\definecolor{warn}{HTML}{B91C1C}

% 계층 색 — Simulink 모델의 블록 색과 같은 규칙을 쓴다
\definecolor{guid}{HTML}{1D4ED8}    % 유도
\definecolor{ctrl}{HTML}{EA580C}    % 제어
\definecolor{thr}{HTML}{D97706}     % 추진기
\definecolor{plant}{HTML}{16A34A}   % 운동모델
\definecolor{nav}{HTML}{7C3AED}     % 항법 · 신호처리
\definecolor{pathc}{HTML}{0D9488}   % 계획 경로 · 항적
\definecolor{errc}{HTML}{C2410C}    % 오차
\definecolor{flow}{HTML}{0369A1}    % 조류 · 바람

% ---- 선과 화살촉 -----------------------------------------------------------
\tikzset{
  sig/.style   = {ink, line width=0.5pt, -{Stealth[length=4pt, width=3pt]}},
  wire/.style  = {ink, line width=0.5pt},
  fb/.style    = {ink, line width=0.5pt, densely dashed,
                  -{Stealth[length=4pt, width=3pt]}},
  asig/.style  = {accent, line width=0.5pt, densely dashed,
                  -{Stealth[length=4pt, width=3pt]}},
  % 블록 — 흰 바탕, 직각, 검은 테두리. 색은 테두리에만 준다
  blk/.style   = {draw=ink, line width=0.55pt, fill=white,
                  minimum width=18mm, minimum height=8.5mm, inner sep=2pt, font=\small},
  gblk/.style  = {blk, draw=guid},
  cblk/.style  = {blk, draw=ctrl},
  tblk/.style  = {blk, draw=thr},
  pblk/.style  = {blk, draw=plant},
  nblk/.style  = {blk, draw=nav},
  % 합산점 · 분기점
  sum/.style   = {draw=ink, line width=0.55pt, fill=white, circle,
                  inner sep=0pt, minimum size=4.6mm, font=\scriptsize},
  dot/.style   = {ink, fill=ink, circle, inner sep=0pt, minimum size=1.5mm},
  % 글자
  sname/.style = {font=\small, inner sep=1.5pt},
  note/.style  = {font=\scriptsize, soft, inner sep=1.5pt},
  wnote/.style = {font=\scriptsize, warn, inner sep=1.5pt},
  anote/.style = {font=\scriptsize, accent, inner sep=1.5pt},
  axis/.style  = {ink, line width=0.5pt},
  tick/.style  = {font=\scriptsize, soft},
  % 묶음 상자 — 서브시스템 표시
  grp/.style   = {draw=soft, line width=0.4pt, densely dotted, rounded corners=1mm,
                  inner sep=3mm},
}
   (assets/src/ 안에서)
% 빌드:     bash _tools/tikz2svg.sh assets/src/w02-node-topic.tex assets/w02-node-topic.svg
%
% 왜 TikZ 인가
%   손으로 SVG 좌표를 쓰면 화살촉·선 굵기·글자 크기가 그림마다 미묘하게 달라진다.
%   그것이 "자료가 어설퍼 보인다"의 정체다. 스타일을 한 번만 정하고 상대좌표로
%   배치하면 좌표를 고쳐도 그림이 무너지지 않는다.
%   수식은 본문과 같은 TeX 글꼴로 나온다.
%
% 한글
%   ko.TeX(DVI 경로)를 쓴다. XeLaTeX 은 TikZ 그래픽을 PDF special 로 내보내는데
%   dvisvgm 이 그것을 받으려면 Ghostscript < 10.01 이 필요하다. 이 PC 에는
%   10.07 이 깔려 있어 그 경로가 막힌다. latex -> DVI -> dvisvgm 이 유일하게
%   동작하는 조합이다.

\usepackage[hangul]{kotex}
\usepackage{tikz}
\usepackage{amsmath}
\usepackage{xcolor}
\usetikzlibrary{arrows.meta, positioning, calc, fit, backgrounds, decorations.pathmorphing}

% ---- 색 --------------------------------------------------------------------
% 강조는 하나만. 나머지는 검정. 흑백 인쇄에서도 읽혀야 한다.
\definecolor{ink}{HTML}{1A1A1A}
\definecolor{soft}{HTML}{666666}
\definecolor{accent}{HTML}{7C3AED}
\definecolor{warn}{HTML}{B91C1C}

% 계층 색 — Simulink 모델의 블록 색과 같은 규칙을 쓴다
\definecolor{guid}{HTML}{1D4ED8}    % 유도
\definecolor{ctrl}{HTML}{EA580C}    % 제어
\definecolor{thr}{HTML}{D97706}     % 추진기
\definecolor{plant}{HTML}{16A34A}   % 운동모델
\definecolor{nav}{HTML}{7C3AED}     % 항법 · 신호처리
\definecolor{pathc}{HTML}{0D9488}   % 계획 경로 · 항적
\definecolor{errc}{HTML}{C2410C}    % 오차
\definecolor{flow}{HTML}{0369A1}    % 조류 · 바람

% ---- 선과 화살촉 -----------------------------------------------------------
\tikzset{
  sig/.style   = {ink, line width=0.5pt, -{Stealth[length=4pt, width=3pt]}},
  wire/.style  = {ink, line width=0.5pt},
  fb/.style    = {ink, line width=0.5pt, densely dashed,
                  -{Stealth[length=4pt, width=3pt]}},
  asig/.style  = {accent, line width=0.5pt, densely dashed,
                  -{Stealth[length=4pt, width=3pt]}},
  % 블록 — 흰 바탕, 직각, 검은 테두리. 색은 테두리에만 준다
  blk/.style   = {draw=ink, line width=0.55pt, fill=white,
                  minimum width=18mm, minimum height=8.5mm, inner sep=2pt, font=\small},
  gblk/.style  = {blk, draw=guid},
  cblk/.style  = {blk, draw=ctrl},
  tblk/.style  = {blk, draw=thr},
  pblk/.style  = {blk, draw=plant},
  nblk/.style  = {blk, draw=nav},
  % 합산점 · 분기점
  sum/.style   = {draw=ink, line width=0.55pt, fill=white, circle,
                  inner sep=0pt, minimum size=4.6mm, font=\scriptsize},
  dot/.style   = {ink, fill=ink, circle, inner sep=0pt, minimum size=1.5mm},
  % 글자
  sname/.style = {font=\small, inner sep=1.5pt},
  note/.style  = {font=\scriptsize, soft, inner sep=1.5pt},
  wnote/.style = {font=\scriptsize, warn, inner sep=1.5pt},
  anote/.style = {font=\scriptsize, accent, inner sep=1.5pt},
  axis/.style  = {ink, line width=0.5pt},
  tick/.style  = {font=\scriptsize, soft},
  % 묶음 상자 — 서브시스템 표시
  grp/.style   = {draw=soft, line width=0.4pt, densely dotted, rounded corners=1mm,
                  inner sep=3mm},
}


\begin{document}
\begin{tikzpicture}[x=1mm, y=1mm]

\def\S{15}   % 1 m

% ===================== 평면도 =====================
\begin{scope}
  \node[font=\small\bfseries, anchor=west] at (-46,30) {평면도 — 위에서 본 모습};

  % 라이다는 한 바퀴를 다 본다 (360 deg) — 10주차 카메라와 가장 다른 점
  \draw[errc, densely dotted, opacity=0.55] (0.700*\S,0) circle (30mm);
  \draw[errc, line width=0.4pt] (31.7,21.2) -- (38,25);
  \node[note, errc, anchor=west] at (38.5,25) {라이다 수평 시야 $360^\circ$};

  % 선체 두 개 (y = +-1.027 m)
  \draw[flow, line width=0.7pt, rounded corners=2mm]
       (-2.374*\S-9,  1.027*\S-3) rectangle (0.9*\S+9,  1.027*\S+3);
  \draw[flow, line width=0.7pt, rounded corners=2mm]
       (-2.374*\S-9, -1.027*\S-3) rectangle (0.9*\S+9, -1.027*\S+3);
  \node[note, flow, anchor=west] at (0.9*\S+11,  1.027*\S) {좌현 $y_b=+1.03$};
  \node[note, flow, anchor=west] at (0.9*\S+11, -1.027*\S) {우현 $y_b=-1.03$};

  % 추진기 셋 — 6주차 배치 그대로
  \draw[thr, line width=0.7pt] (-2.374*\S-4,  1.027*\S-3) rectangle (-2.374*\S+4,  1.027*\S+3);
  \draw[thr, line width=0.7pt] (-2.374*\S-4, -1.027*\S-3) rectangle (-2.374*\S+4, -1.027*\S+3);
  \node[note, thr, anchor=west] at (-2.374*\S-9, 1.027*\S+6.5) {좌·우 추진기 $x_b=-2.37$};

  \draw[thr, line width=0.7pt] (0.750*\S-3,-4) rectangle (0.750*\S+3,4);
  \draw[thr, -{Stealth[length=4pt,width=3pt]}] (0.750*\S,4) -- (0.750*\S,11);
  \node[note, thr, anchor=west] at (0.750*\S+2,11)
       {선수추진기 $x_b=0.75$ — 옆으로 민다};

  % 기준점과 축
  \fill[ink] (0,0) circle (1.1);
  \node[note, anchor=north east] at (-3,-1.8) {\texttt{base\_link}};
  \draw[axis, -{Stealth[length=4pt,width=3pt]}] (0,0) -- (7,0);
  \node[note, anchor=south west] at (1.5,0.8) {$x_b$};
  \draw[axis, -{Stealth[length=4pt,width=3pt]}] (0,0) -- (0,11);
  \node[note, anchor=east] at (-1,9.5) {$y_b$ 좌현};

  % 라이다 — 선수추진기 바로 위라 평면도에서는 겹쳐 보인다
  \draw[errc, line width=0.9pt] (0.700*\S,0) circle (2.6);
  \fill[errc] (0.700*\S,0) circle (0.9);
  \draw[errc, line width=0.4pt] (0.700*\S-1.8,-1.9) -- (-6,-10);
  \node[note, errc, anchor=east] at (-6.5,-10)
       {3D LiDAR $x_b=0.70$ — 선수추진기 바로 위};

  % GPS / IMU
  \fill[flow] (-0.850*\S,0) circle (1.4);
  \node[note, flow, anchor=south] at (-0.850*\S,2) {GPS};
  \fill[nav]  (0,-0.200*\S) circle (1.2);
  \node[note, nav, anchor=north west] at (1.5,-4.2) {IMU};
\end{scope}

% ===================== 측면도 =====================
\begin{scope}[shift={(0,-70)}]
  \node[font=\small\bfseries, anchor=west] at (-46,36) {측면도 — 높이를 정한 근거};

  \draw[flow, densely dashed] (-44,0) -- (88,0);
  \node[note, flow, anchor=south] at (70,0.8) {수면};
  \draw[flow, line width=0.7pt, rounded corners=2mm]
       (-2.374*\S-9,-3) rectangle (0.9*\S+9,3);

  \fill[ink] (0,0) circle (1.1);
  \node[note, anchor=north] at (0,-4.5) {\texttt{base\_link}};

  % 전장함 — 이것이 범인이다
  \fill[warn, opacity=0.18]
       (-0.295*\S,1.295*\S) rectangle (0.295*\S,1.765*\S);
  \draw[warn, line width=0.6pt]
       (-0.295*\S,1.295*\S) rectangle (0.295*\S,1.765*\S);
  \node[note, warn, anchor=east] at (-6,22.9) {\texttt{cpu\_cases\_link}};

  % 높이 셋을 오른쪽 한 줄로 읽는다
  \draw[errc, densely dotted] (13,1.900*\S) -- (40,1.900*\S);
  \draw[errc, line width=0.4pt] (40,1.900*\S) -- (42,30.5);
  \node[note, errc, anchor=west] at (42.5,30.5)
       {3D LiDAR $z_b=1.90$ — 고른 자리, 자기 배 $0\%$};

  \draw[warn, densely dotted] (19,1.765*\S) -- (40,1.765*\S);
  \draw[warn, line width=0.4pt] (40,1.765*\S) -- (42,26.5);
  \node[note, warn, anchor=west] at (42.5,26.5) {전장함 윗면 $z_b=1.765$};

  % 기둥과 라이다
  \draw[soft, line width=0.5pt] (0.700*\S,0) -- (0.700*\S,1.900*\S);
  \fill[errc] (0.700*\S,1.900*\S) circle (1.8);

  % 틀렸던 자리
  \draw[soft, line width=0.7pt] (0.700*\S,1.600*\S) circle (1.3);
  \draw[soft, densely dotted] (13,1.600*\S) -- (40,1.600*\S);
  \draw[soft, line width=0.4pt] (40,1.600*\S) -- (42,22);
  \node[note, soft, anchor=west] at (42.5,22)
       {$1.60$ — 계산으로 골랐던 자리, $82\%$ 가 자기 배};

  % 여유
  \draw[axis, {Stealth[length=3pt,width=2.5pt]}-{Stealth[length=3pt,width=2.5pt]}]
       (6.5,1.765*\S) -- (6.5,1.900*\S);
  \draw[axis, line width=0.4pt] (6.5,1.900*\S) -- (6.5,29.0);
  \node[note, anchor=south east] at (7,29.0) {여유 $13.5$ cm};

  % 세로 시야 아래쪽 광선 -15 deg
  \draw[pathc, line width=0.7pt] (0.700*\S,1.900*\S) -- (38,21.13);
  \draw[pathc, line width=0.5pt] (0.700*\S+7,1.900*\S) arc (0:-15:7);

  % 그 광선이 수면에 닿는 곳 — 치수선을 끊어 그린다
  \draw[pathc, densely dotted] (0.700*\S,-3) -- (0.700*\S,-12.5);
  \draw[pathc, densely dotted] (80,0) -- (80,-12.5);
  \fill[pathc] (80,0) circle (1.2);
  \draw[pathc, {Stealth[length=3pt,width=2.5pt]}-{Stealth[length=3pt,width=2.5pt]}]
       (0.700*\S,-13) -- (80,-13);
  \draw[pathc, line width=0.6pt] (46,-14.5) -- (47.5,-11.5);
  \draw[pathc, line width=0.6pt] (48.5,-14.5) -- (50,-11.5);
  \node[note, pathc, anchor=north, text width=70mm, align=center] at (45,-14.8)
       {가장 아래 광선($-15^\circ$)이 수면에 닿는 곳 $7.09$ m\\
        — 이보다 가까운 물 위는 라이다가 못 본다};

  % 선수추진기 프로펠러
  \fill[thr] (0.750*\S,0.318*\S-0.509*\S) circle (1.2);
  \node[note, thr, anchor=north west] at (0.750*\S+2,-3.4) {선수추진기};
\end{scope}

% ===================== 주의 =====================
\node[wnote, anchor=north west, text width=130mm] at (-46,-96)
     {높이를 \textbf{계산으로 1.60 m} 라 정했다가 실측에서 틀린 것을 알았다 — 들어온 값의
      82 \% 가 장애물이 아니라 자기 배였다. $1.75 \to 1.80$ 사이에서 82 \% 가 0 \% 로 뚝 떨어진다.
      파도에 기우는 것(여유 $13.5$ cm $\approx$ 횡동요 $18^\circ$)까지 보아 \textbf{1.90 m} 로 둔다.
      값은 비싸지 않다 — 부표를 처음 보는 거리가 $4.9 \to 5.2$ m 로 느는 것뿐이다.
      이 그림의 $y_b$ 는 \textbf{좌현이 양(+)} 이다 — 5주차 이후 NED 수식과 부호가 반대다};

\end{tikzpicture}
\end{document}
